\documentclass[a4paper]{article}
\usepackage{amsfonts}
\usepackage{amsmath}
\usepackage{amssymb}
\usepackage{graphicx}     

\begin{document}
  
\noindent \large Auteur: Stan Van Campenhout \\
\noindent \large Studierichting:  Informatica\\
\noindent \large Oefeningenreeks 4A nr. 44.14\\

\medskip

\normalsize

$ \displaystyle\int\dfrac{2x\left(x+5\right)}{x^4-1}dx = \displaystyle\int\dfrac{2x^2+10x}{\left(x^2-1\right)\left(x^2+1\right)}dx
= \displaystyle\int\dfrac{2x^2+10x}{\left(x-1\right)\left(x+1\right)\left(x^2+1\right)}dx $\\

$ \rightarrow \dfrac{2x^2+10x}{\left(x-1\right)\left(x+1\right)\left(x^2+1\right)} = \dfrac{A}{x-1} + \dfrac{B}{x+1} + \dfrac{Cx + D}{x^2 +1} $\\

$ \Rightarrow A = \dfrac{2 + 10}{\left(1+1\right)\left(1+1\right)} = \dfrac{12}{4} = 3 $\\

$ \Rightarrow B = \dfrac{2 - 10}{\left(-1-1\right)\left(1+1\right)} = \dfrac{-8}{-4} = 2 $\\

$ \Rightarrow C\left(i\right) + D = \dfrac{-2+10i}{\left(i-1\right)\left(i+1\right)} = \dfrac{10i - 2}{-2} = 1 - 5i $\\

$ \Rightarrow C = -5, D = 1 $\\

$ = 3\displaystyle\int\dfrac{1}{x-1}dx + 2\displaystyle\int\dfrac{1}{x+1} + \displaystyle\int\dfrac{-5x + 1}{x^2 + 1}dx $\\

$ = 3\ln|x-1| + 2\ln|x+1| - \dfrac{5}{2}\displaystyle\int\dfrac{1}{x^2+1}d\left(x^2 + 1\right) + \displaystyle\int\dfrac{1}{x^2+1}dx $\\

$ = 3\ln|x-1| + 2\ln|x+1| - \dfrac{5}{2}\ln|x^2+1| + \operatorname{Bgtan}x + c $\\

\begin{thebibliography}{999}
%\bibitem{a} Referentie
\end{thebibliography}
\end{document} 
